\documentclass[11pt]{article}
\usepackage[a4paper,margin=28mm]{geometry}
\usepackage[T1]{fontenc}
\usepackage{lmodern}
\usepackage{amsmath,amssymb,amsthm}
\usepackage{microtype}
\usepackage[colorlinks=true,linkcolor=blue,citecolor=blue,urlcolor=blue]{hyperref}
\hypersetup{pdftitle={Triangle-free circulants with odd degree equal to their independence number},
pdfauthor={},pdfsubject={A constructive proof of an existence statement for circulant graphs},
pdfkeywords={circulant graph, triangle-free graph, independence number}}
\newtheorem{theorem}{Theorem}
\newtheorem{lemma}[theorem]{Lemma}
\newcommand{\Z}{\mathbb Z}
\newcommand{\Cay}{\operatorname{Cay}}
\title{Triangle-free circulants with odd degree\\equal to their independence number}
\author{}
\date{October 13, 2026}

\begin{document}
\maketitle
\begin{abstract}
For every integer \(k\ge1\) and every even integer \(n\) satisfying
\(6k+2\le n\le8k\), we construct a triangle-free circulant graph on \(n\)
vertices whose degree and independence number both equal \(2k+1\).
This proves the existence assertion in Conjecture~3.8 of
Ayat, Ayat and Ghahramani (2020).
For \(4\mid n\), the construction extends their upper-endpoint construction,
and an upper bound on independent sets follows from circular gaps.
For \(n\equiv2\pmod4\), a Chinese remainder isomorphism gives a two-layer
description, and expansion by a cyclic interval gives the upper bound.
The proof is elementary and does not depend on finite computation.
No claim of bibliographic priority is made.
\end{abstract}

\section{Statement and relation to earlier constructions}

Write \(\alpha(G)\) for the independence number of a finite simple graph \(G\).
For \(S=-S\subseteq\Z_n\setminus\{0\}\), the circulant
\(\Cay(\Z_n,S)\) has vertex set \(\Z_n\), with \(x\) adjacent to \(y\)
if and only if \(y-x\in S\).
Ayat, Ayat and Ghahramani call a triangle-free circulant whose independence
number equals its degree an \(A\)-graph~\cite{Ayat2020}.
Their Conjecture~3.8 is the following assertion.

\begin{theorem}\label{thm:main}
For every integer \(k\ge1\) and every even integer \(n\) with
\[
6k+2\le n\le8k,
\]
there is a simple undirected triangle-free circulant graph on \(n\) vertices
that is \((2k+1)\)-regular and has independence number \(2k+1\).
\end{theorem}

Theorems~3.6 and~3.7 of~\cite{Ayat2020} establish the endpoints \(n=8k\)
and \(n=6k+2\), respectively.
Our construction for \(4\mid n\) retains the short odd jumps and the
circular-gap argument of their \(n=8k\) construction, while replacing
the single antipodal jump by an interval of even jumps.
The lower endpoint also belongs to the classical Andr\'asfai family,
whose order at degree \(d\) is \(3d-1\); see
Theorem~11 of~\cite{Bickle2027}.
Our proof treats every intermediate even \(n\), as well as both endpoints.

Sidorenko's construction has even degree~\cite{Sidorenko1991};
its exact circulant form is also stated in Theorem~3.5
of~\cite{Ayat2020}.
Brandt's Lemma~3.2~\cite{Brandt2010} realizes every rational
order-to-degree ratio in \([3,4)\) by a Sidorenko graph.
Specifying a ratio does not specify the order and the degree separately:
for a ratio \(b/a\), that construction uses order \(4b\) and degree \(4a\).
Thus it does not by itself give the prescribed odd degree in
Theorem~\ref{thm:main}.

A later paper of Bickle, Cream and Gould~\cite{Bickle2027},
available online on July 21, 2026, restates the even-degree family,
the Andr\'asfai family, and the odd-degree upper endpoint
(Theorems~10--12). Its equal-size blow-ups preserve the relevant
properties and circulant structure (Theorem~13 and Lemma~14).
These results do not directly fill the whole prescribed parameter range.
For example, at degree \(7\), the two endpoint families give orders \(20\)
and \(24\), whereas Theorem~\ref{thm:main} also requires order \(22\).
An even-degree graph still has even degree after such a blow-up.
A nontrivial blow-up yielding degree \(7\) must start with degree \(1\)
and use multiplicity \(7\); the relevant connected base graph is \(K_2\),
giving order \(14\).
This comparison concerns the coverage of these particular results;
it is not a proof that no other earlier construction exists.
The literature search was not exhaustive, and no priority claim is made.

\section{Elementary tools}

All sums of subsets below are taken in the indicated finite abelian group.
A cyclic interval of length \(\ell\le m\) in \(\Z_m\) is a translate
of \(\{0,1,\ldots,\ell-1\}\).

\begin{lemma}\label{lem:triangle}
Let \(S=-S\subseteq\Gamma\setminus\{0\}\), where \(\Gamma\) is a finite
abelian group. The graph \(\Cay(\Gamma,S)\) is triangle-free if and only if
\((S+S)\cap S=\varnothing\).
Every nonempty triangle-free \(d\)-regular graph has independence number
at least \(d\).
\end{lemma}
\begin{proof}
If \(s,t,s+t\in S\), then \(0,s,s+t\) are three distinct vertices forming
a triangle. Conversely, translate one vertex of a triangle to \(0\).
If its other vertices are \(x,y\), then \(x,y,y-x\in S\), and
\(y=x+(y-x)\in(S+S)\cap S\).
For the last assertion, the \(d\) neighbors of any vertex form an
independent set.
\end{proof}

\begin{lemma}\label{lem:interval}
If \(\varnothing\ne A\subseteq\Z_m\) and \(J\) is a cyclic interval
of length \(1\le\ell\le m\), then
\[
|A+J|\ge\min\{m,|A|+\ell-1\}.
\]
\end{lemma}
\begin{proof}
Translate \(J\) to \(\{0,\ldots,\ell-1\}\), without changing the
cardinality of the sumset. Every nonempty proper subset \(B\) of \(\Z_m\)
satisfies \(|B\cup(B+1)|\ge|B|+1\).
Otherwise finiteness gives \(B+1=B\), and invariance under the generator
\(1\) forces \(B=\Z_m\).
Starting with \(B_0=A\), set \(B_j=A+\{0,\ldots,j\}\).
Since \(B_{j+1}=B_j\cup(B_j+1)\), the size increases by at least one
at each step until it reaches \(m\). This proves the bound.
No primality assumption on \(m\) is needed.
\end{proof}

We also use the following edge count. If \(H\) is \(d\)-regular with \(d>0\)
and \(I\) is independent, then
\[
d|I|=e(I,V(H)\setminus I)\le d|V(H)\setminus I|,
\qquad\text{so}\qquad \alpha(H)\le |V(H)|/2.
\tag{1}\label{eq:regular}
\]

\section{\texorpdfstring{The case \(n\equiv0\pmod4\)}{The case n congruent to 0 modulo 4}}

Set
\[
r=\frac{8k-n}{4},\qquad a=k-r,\qquad m=\frac n2.
\]
Then \(k=a+r\), \(n=8a+4r\), and the parameter conditions are equivalent to
\[
a\ge1,\qquad 0\le r\le a-1.
\tag{2}\label{eq:firstparameters}
\]
In \(\Z_n\), define
\[
L=\{\pm1,\pm3,\ldots,\pm(2a-1)\},\qquad
E=\{m+2j:-r\le j\le r\},\qquad S=L\cup E.
\tag{3}\label{eq:firstset}
\]
The elements of \(L\) are odd, whereas those of \(E\) are even.
The latter have distinct integer representatives between
\(m-2r=4a\) and \(m+2r=n-4a\).
Both sets are symmetric and exclude \(0\), and neither has repeated
residues in its defining list.
Consequently \(G=\Cay(\Z_n,S)\) is a simple undirected circulant of degree
\[
|S|=2a+(2r+1)=2k+1.
\]

\subsection*{Triangle-freeness}

Every element of \(E\) has circular distance at least \(4a\) from \(0\).
Each element of \(L+L\) has an integer representative in
\([-4a+2,4a-2]\), so \((L+L)\cap E=\varnothing\).
Since \(2m=n\), each element of \(E+E\) has a representative in
\([-4r,4r]\). By~\eqref{eq:firstparameters}, \(4r<4a\), and hence
\((E+E)\cap E=\varnothing\).
The directed edge differences around a triangle must sum to zero.
Parity allows only two \(L\)-differences and one \(E\)-difference,
or three \(E\)-differences. Symmetry reduces these possibilities to the
two intersections just excluded. In particular, an equality
\(l+e=l'\), with \(l,l'\in L\) and \(e\in E\), would give
\(e=l'-l\in L+L\).
Thus \(G\) is triangle-free.

\subsection*{Independence number}

Let \(I\) be an independent set and let \(s=|I|\).
If all its vertices have the same parity, note that \(m\) is even and
\(m\in S\). The antipodal edges \(\{x,x+m\}\) partition that parity class
into \(n/4\) pairs. Hence
\[
s\le n/4=2a+r\le2a+2r+1=2k+1.
\]
The empty independent set satisfies this bound as well.

Otherwise order the vertices of \(I\) around the circle and consider the
\(s\) positive gaps, including the gap from the last vertex back to the
first. Their sum is \(n\). The number \(q\) of odd gaps is at least two,
because a circular sequence containing both parities has at least two
parity changes. Each even gap is at least \(2\), and each odd gap is at
least \(2a+1\), since the smaller positive odd differences belong to \(L\).
As \(a\ge1\), this gives
\[
n\ge2(s-q)+(2a+1)q
 =2s+(2a-1)q\ge2s+4a-2.
\]
It follows that
\[
s\le\frac{n-4a+2}{2}=2a+2r+1=2k+1.
\]
The opposite inequality follows from Lemma~\ref{lem:triangle}.
Therefore \(\alpha(G)=2k+1\).

\section{\texorpdfstring{The case \(n\equiv2\pmod4\)}{The case n congruent to 2 modulo 4}}

Set
\[
m=\frac n2,\qquad t=\frac{n-4k+2}{4}.
\]
Then \(m\) is odd, \(t\) is an integer, and
\[
m=2k+2t-1,\qquad t\le k\le2t-2.
\tag{4}\label{eq:secondparameters}
\]
Indeed, the lower bound on \(n\) gives \(k\le2t-2\), while this residue
class has \(n\le8k-2\), giving \(t\le k\).
In particular \(t\ge2\), so this case has no parameter with \(k=1\).

Since \(m\) is odd, the Chinese remainder theorem gives a group isomorphism
\[
\phi:\Z_{2m}\longrightarrow\Z_m\times\Z_2,\qquad
x\longmapsto(x\bmod m,x\bmod2).
\]
In \(\Z_m\), put
\[
T=\{-(t-1),\ldots,t-1\},\qquad U=\{2t-1,2t,\ldots,2k\},
\]
and define
\[
Q=(T\times\{1\})\cup(U\times\{0\}).
\tag{5}\label{eq:secondset}
\]
Both intervals have distinct residues. We have \(T=-T\) and \(m-U=U\);
thus \(Q=-Q\), and \(Q\) does not contain the identity \((0,0)\).
Although \(0\in T\), its second coordinate is \(1\), so it creates no loop.
Moreover,
\[
|Q|=(2t-1)+(2k-2t+2)=2k+1.
\]
The graph \(G=\Cay(\Z_m\times\Z_2,Q)\) is therefore a simple undirected
\((2k+1)\)-regular circulant through \(\phi\).
Equivalently, its connection set on \(\Z_n\) is
\[
S=\{x\in\{1,\ldots,n-1\}:
 (x\text{ odd and }x\bmod m\in T)
 \text{ or }(x\text{ even and }x\bmod m\in U)\}.
\tag{6}\label{eq:directcrt}
\]

\subsection*{Triangle-freeness}

The second coordinates of the three directed edge differences around a
triangle sum to \(0\) in \(\Z_2\). A triangle therefore has either two
cross-layer edges or none.

For two cross-layer edges, it is enough, by symmetry, to check
\((T+T)\cap U=\varnothing\).
Each element of \(T+T\) has an integer representative in
\([-2t+2,2t-2]\).
Its nonnegative representatives are below \(2t-1\), and its negative
representatives become at least \(m-(2t-2)=2k+1\) after reduction modulo
\(m\). Both parts avoid \(U\).

For three edges within a layer, the integer sums from \(U+U\) lie in
\([4t-2,4k]\).
If a sum has not wrapped around modulo \(m\), it exceeds \(2k\), since
\(4t-2>2k\) by~\eqref{eq:secondparameters}.
After one wrap it is at most
\[
4k-m=2k-2t+1\le2t-3<2t-1.
\]
There cannot be two wraps, because \(4k<2m\).
Thus \((U+U)\cap U=\varnothing\).
The graph is triangle-free and has independence number at least \(2k+1\).

\subsection*{Independence number}

Write an arbitrary independent set as
\[
(A\times\{0\})\cup(B\times\{1\}),\qquad A,B\subseteq\Z_m.
\]
If both \(A\) and \(B\) are nonempty, the absence of cross-layer edges gives
\(B\cap(A+T)=\varnothing\).
By Lemma~\ref{lem:interval},
\[
|A+T|\ge\min\{m,|A|+2t-2\}.
\]
Because \(B\ne\varnothing\), the sumset is not all of \(\Z_m\).
Consequently \(|A|+2t-2<m\), and
\[
|A|+|B|\le |A|+m-|A+T|\le m-2t+2=2k+1.
\]

If one layer is empty, the independent set lies in
\(H=\Cay(\Z_m,U)\), whose degree is
\(2(k-t+1)\ge2\).
Applying~\eqref{eq:regular} and using \(m\) odd yields
\[
|A|+|B|\le\lfloor m/2\rfloor=k+t-1\le2k-1<2k+1.
\]
This also includes the empty independent set.
Combining the bounds proves \(\alpha(G)=2k+1\).

\section{Endpoints, verification and status}

Every even \(n\) belongs to exactly one of the two cases, which proves
Theorem~\ref{thm:main}. For clarity, the boundary parameters are as follows.
When \(k=1\), the only allowed order is \(8\), and the first construction
has \(a=1,r=0,S=\{1,4,7\}\).
At \(n=8k\), the first construction has \(a=k,r=0\).
At \(n=6k+2\), odd \(k\) gives \(a=(k+1)/2,r=(k-1)/2=a-1\),
whereas even \(k\) gives \(t=(k+2)/2\) and \(k=2t-2\).
The largest order in the second residue class is \(8k-2\), for \(k\ge2\);
here \(t=k\), and the within-layer degree remains \(2\).
All strict inequalities used in the proof remain valid at these endpoints.

The accompanying standard-library Python program implements the two
displayed connection sets directly and computes exact independence
numbers by the deletion recurrence
\[
\alpha(G[P])=\max\bigl\{\alpha(G[P\setminus\{v\}]),
  1+\alpha(G[P\setminus N[v]])\bigr\}.
\]
It does not use the asserted bound \(2k+1\) to prune the search.
The supplied reproducibility files record a run over all \(21\) allowed
pairs with \(1\le k\le6\).
Such finite checks test the implementation; they do not establish the
unbounded theorem, whose proof is given above.

This work was prepared with AI assistance.
It has not yet been reviewed by a human expert or formalized in a proof
assistant. The comparison with the cited literature does not establish
bibliographic priority.

\begin{thebibliography}{9}
\raggedright
\bibitem{Ayat2020}
S.~Masih Ayat, S.~Majid Ayat and M.~Ghahramani,
The independence number of circulant triangle-free graphs,
\emph{AIMS Mathematics} \textbf{5}(4) (2020), 3741--3750.
\href{https://doi.org/10.3934/math.2020242}{doi:10.3934/math.2020242}.

\bibitem{Sidorenko1991}
A.~F.~Sidorenko,
Triangle-free regular graphs,
\emph{Discrete Mathematics} \textbf{91}(2) (1991), 215--217.
\href{https://www.sciencedirect.com/science/article/pii/0012365X9190114H}{Publisher record}.

\bibitem{Brandt2010}
S.~Brandt,
Triangle-free graphs whose independence number equals the degree,
\emph{Discrete Mathematics} \textbf{310} (2010), 662--669.
\href{https://doi.org/10.1016/j.disc.2009.05.021}{doi:10.1016/j.disc.2009.05.021}.

\bibitem{Bickle2027}
A.~Bickle, M.~Cream and R.~J.~Gould,
Chorded cycles in Chv\'atal--Erd\H{o}s graphs,
\emph{Discrete Mathematics} \textbf{350}(1) (2027), article 115336.
Available online July 21, 2026.
\href{https://doi.org/10.1016/j.disc.2026.115336}{doi:10.1016/j.disc.2026.115336}.
\end{thebibliography}
\end{document}

